\section{Introduction}
How many positive integers \(n\le X\) have the property that both
\(n\) and \(n+1\) are sums of two squares? Write \(b_d(n)\) for
the indicator of the positive values of \(x^2+dy^2\), with
\(x,y\in\Z\), and put \(b_d(n)=0\) for \(n\le0\).
Coordinates are unrestricted and may vanish; each numerical value
is counted once. Under the analytic inputs stated below, we obtain
\begin{equation}\label{eq:ck-headline}
 \#\{1\le n\le X:b_1(n)=b_1(n+1)=1\}
                         \sim\frac{X}{2\log X}.
\end{equation}
This is the shift-one Connors--Keating prediction. It belongs to
Smilansky's five-form family of pair-correlation predictions.

\begin{corollary}\label{cor:smilansky-five-pairs}
Under the inputs of Section~\ref{sec:inputs}, for
\(d\in\{1,2,3,4,7\}\) and fixed \(h\ne0\),
\[
 \sum_{1\le n\le X}b_d(n)b_d(n+h)
       =\mathscr T_{d,h}\frac X{\log X}+o_{d,h}(X/\log X),
\]
with precisely the coefficients predicted in
\cite[Section 6.3, equations (6.2)--(6.4)]{smilansky}.
In particular,
\begin{equation}\label{eq:minus-eight-shift-one}
 \sum_{n\le X}b_2(n)b_2(n+1)
                       \sim\frac1{\sqrt2}\frac X{\log X}.
\end{equation}
\end{corollary}

The coefficients and their exceptional-prime conditions are evaluated
in Section~\ref{sec:quadratic-pairs}. In particular these are ordinary
value counts, with no weights for the number of coordinate representations.

\subsection{A second headline: genus governs primitive value counts}

Fix one negative fundamental discriminant \(D\) and two fixed
discriminants \(\Delta_i=f_i^2D\). Let \(Q_i,Q'_i\) be primitive
positive definite integral binary forms of discriminant \(\Delta_i\),
in the same genus for each \(i=1,2\). Thus both slots belong to
orders in the one quadratic field \(\Q(\sqrt D)\); no mixed-field
assertion is made. If \(P_Q\) is the set of positive values of
\(Q\) at coprime coordinates and
\(R_{Q_1,Q_2}(N)=\#\{(a,b)\in P_{Q_1}\times P_{Q_2}:a+b=N\}\),
then the classical route gives
\begin{equation}\label{eq:genus-headline}
 R_{Q_1,Q_2}(N)=R_{Q'_1,Q'_2}(N)
                  \left(1+O((\log N)^{-1/5})\right)
\end{equation}
on sufficiently large common locally admissible targets.
The two slots may have different genera and conductors; replacement
is within each respective discriminant and genus. Under the recent
analytic inputs we also obtain an explicit mixed leading term of
order \(NJ_S(N)/\log N\), with relative error \(o(1)\).
There is no class-number divisor: a numerical value can be represented
by every class in its compatible genus. Section~\ref{sec:forms}
states the general results before giving the ideal-class and local-filter
proofs; Section~\ref{sec:norms} recovers the principal specialization.

The counting conventions and strengths are different:
\begin{center}\small
\begin{tabular}{@{}>{\raggedright\arraybackslash}p{.19\linewidth}>{\raggedright\arraybackslash}p{.34\linewidth}>{\raggedright\arraybackslash}p{.41\linewidth}@{}}
\hline
Result family&What is counted&Conclusion\\\hline
CK/Smilansky&Ordinary values; unrestricted coordinates&Fixed-shift asymptotic with the stated little-\(o\) error\\[1mm]
Subgroup/coset sums&Integers with prescribed prime support&Target-uniform asymptotic with a logarithmic saving\\[1mm]
Primitive-form sums&Values admitting coprime coordinates in two fixed orders of one field&Explicit leading term with relative \(o(1)\)\\[1mm]
Genus replacement&Comparison of two primitive-form value counts&Relative logarithmic saving through the classical route\\\hline
\end{tabular}
\end{center}
The first three asymptotic families use the imports in
Section~\ref{sec:inputs}; genus replacement does not.
A \emph{primitive form} \([a,b,c]\) has \(\gcd(a,b,c)=1\).
A \emph{primitive representation} \(Q(x,y)\) has
\(\gcd(x,y)=1\). \emph{Coprime summands} in \(a+b=N\)
means \(\gcd(a,b)=1\), a separate condition imposed only when stated.
For example, \(4=2^2+0^2\) is an ordinary sum of two squares,
but any representation of four has both coordinates even and is
not primitive. Consequently the primitive-form theorem does not
extend the ordinary CK statement to arbitrary forms. Its class-error
saving also does not supply a logarithmic remainder for the full
conductor-prefix leading term.

\subsection{The common mechanism, in the CK example}

The central result is the indicator-to-model comparison of
Theorem~\ref{thm:model-comparison}. Multiplication by suitable primes
preserves an indicator and its arithmetic model. One-variable means
center their difference, and a positive sparse sieve controls exceptional
prime-factor records. A shared graph reduction then leaves two frequency
estimates: one for fixed shifts and one for reflected targets. Their
assembly replaces the indicators in the binary count by the models.
The model correlation can then be evaluated from local arithmetic.

Here is the model for \(b_1\). At outer scale \(X\), put
\(L=\log X\), \(y=\exp(L^{2/5})\),
\(\mathcal P_y=\{p:p\le y\}\), and
\(V_y=\prod_{p\le y}(1-1/p)\). Write every \(n\) uniquely as
\(n=rv\), where \(r=n_{\mathcal P_y}\) contains its full small-prime
powers and \(v\) has no prime factor at most \(y\).
The rough ordinary-value series is
\[
 F_y(s)=\prod_{\substack{p>y\\p\equiv1\ (4)}}(1-p^{-s})^{-1}
         \prod_{\substack{p>y\\p\equiv3\ (4)}}(1-p^{-2s})^{-1}
       =(s-1)^{-1/2}H_y(s).
\]
For a sufficiently deep fixed \(K\), its smooth density expansion is
\[
 M(v)=\sum_{j=0}^K\frac{H_y^{(j)}(1)}{j!}
            \frac{(\log v)^{-1/2-j}}{\Gamma(1/2-j)},\qquad
 \mathcal D(rv)=b_1(r)\bigl(1+\chi_{-4}(v)\bigr)\frac{M(v)}{V_y}.
\]
These formulas apply on the polynomial rough scales used in the proof;
the smooth bounded extension is specified in
Section~\ref{sec:model-preparation} and the quadratic application.
The leading term of \(M\) is \(H_y(1)/\sqrt{\pi\log v}\).
Dividing by \(V_y\), the density of rough integers, gives the
conditional density on the rough variable. Every allowed rough value
has \(\chi_{-4}(v)=1\). The correction \(1+\chi_{-4}(v)\)
retains this bias; average density alone would miss a leading contribution.
For a split prime \(p\le y\), multiplication by \(p\) changes only
\(r\), preserves \(b_1(r)\), and leaves both \(v\) and
\(\mathcal D\) unchanged. Replacing
\(b_1(n)b_1(n+h)\) by \(\mathcal D(n)\mathcal D(n+h)\)
with an error \(o(X/\log X)\) is the analytic task. The remaining
model product gives the coupled local factors and the coefficient
\(1/2\) in \eqref{eq:ck-headline}.

\subsection{Semigroups and the reflected-target obstacle}

The same mechanism reaches a different geometry. Let
\(C_{12}=\{n\ge1:p\mid n\Rightarrow p\equiv1\pmod{12}\}\),
including one. Every sufficiently large \(N\equiv2\pmod{12}\)
is a sum of two elements of \(C_{12}\), and
the ordered value-pair count
\(R_{12}(N)=\#\{(a,b)\in C_{12}^2:a+b=N\}\) satisfies
\begin{equation}\label{eq:opening}
 R_{12}(N)=A_{12}J_{12}(N)\frac{N}{(\log N)^{3/2}}
                  \left(1+O((\log N)^{-\delta})\right),
 \qquad\delta>0.
\end{equation}
The positive factors are explicit in Theorem~\ref{thm:c12}.
For \(C^-_4=\{n\ge1:p\mid n\Rightarrow p\equiv3\pmod4\}\),
every sufficiently large even target likewise has an explicit
asymptotic of order \(N/\log N\), with a positive logarithmic saving.
The subgroup and coset theorems allow fixed arithmetic data and retain
all surviving characters, including the complex branches forced by
the prime class \(3\pmod{10}\). Numerical saving exponents are
collected in Section~\ref{sec:assembly}.

One-variable density does not control the particular line \(a+b=N\),
and averages can conceal exceptional targets. The nonconstant linear
parts of \(n\) and \(N-n\) are dependent, so the independent-linear-form
correlation theorems of Loughran--Matthiesen do not directly apply
\cite[Theorems 1.9 and 3.10]{loughranmatthiesen}.
The reflected frequency argument must retain cancellation and
normalization when the target enters the rational denominators.

Iwaniec--Pomyka\l a's quartic-norm work is a direct antecedent
\cite{iwaniecpomykala}. For \(\Q(\zeta_{12})\), its splitting conditions
select the primes \(1\pmod{12}\). Their Theorem~2 gives conditional
pointwise order bounds for coprime sums; hypothesis (P4) is uniform
in growing shifts. Their Theorem~1 treats fixed differences
unconditionally. They count numerical values once.
Our unrestricted leading term gives a coprime asymptotic by M\"obius
inversion, with the extra target factor in
Corollary~\ref{cor:coprime-c12}. Historical priority beyond these
precise comparisons remains separate from the proof.

\subsection{Attribution, dependencies and scope}

Walsh's multiplicative-invariance framework treats positive integer
divisor orders and a fixed nonzero shift \cite{walsh}.
He explicitly proposes sums-of-two-squares indicator correlations
in Section~1.2.4, pp.~5--6. The application idea is already present
there; the work here supplies the complete indicator models,
arithmetic hypothesis checks and local factors, together with the
uniform reflected-target argument. Exact common-prime invariance,
centered one-variable rough means and a positive sparse envelope are
the independently checkable assumptions of the comparison theorem.
No centered binary cancellation estimate is assumed.

The recent imports are the weighted-dilation root-moment theorem
\cite[Section 3]{weightedgraphs} and the all-conductor Dirichlet
zero-free half-plane \(\Re s>11/12\)
\cite[Part I, Theorem 3.1]{zerofree}. The stronger Part II claim is
not needed. Walsh states that classical estimates could recover his
fixed-shift divisor theorem \cite[p.~2]{walsh}, while his printed
Lemma~5.1 and Section~6.4 still use the recent zero-free input.
Our fixed-shift moduli are polynomial-logarithmic, but their Mellin
heights reach \(x(\log x)^C\). Section~\ref{sec:fixed-classical-boundary}
proves a classical intermediate-height estimate and isolates the
remaining high-frequency obligation. Both recent inputs remain in
the fixed-shift theorem. A future replacement there would still leave
the reflected proof's growing-conductor issue and the graph input.

For primitive values in the fixed field \(\Q(\sqrt D)\), the classical
Kloosterman, weighted-correlation and positive-restoration route gives
pointwise order bounds. The class-error sieve then proves
\eqref{eq:genus-headline} independently of the recent inputs.
Its full classical core argument is Appendix~\ref{sec:older-route}.
The common mixed local-filter calculation gives the explicit leading
coefficient under the recent binary theorem; the principal norm case
is a specialization. Constants and thresholds depend on the fixed
forms and orders, and no numerical onset or complete small-exception
classification is asserted here. The exact classifications and
certificates belong to the companion article, and the deposited
window paper addresses where summands may be placed.

\subsection{Notation}

\begin{center}\small
\begin{tabular}{@{}>{\raggedright\arraybackslash}p{.28\linewidth}>{\raggedright\arraybackslash}p{.66\linewidth}@{}}
\hline
Symbol&Role\\\hline
\(X\), \(L=\log X\), \(h\)&Outer scale, its logarithm, and the fixed shift\\
\(\mathcal P_y\), \(n_{\mathcal P_y}\), \(V_y\)&Small primes, complete small-prime part, and rough-integer density\\
\(f_i,\mathcal D_i,\xi_i\)&Indicators, arithmetic models, and centered errors\\
\(D,D_i\), \(U_i\)&Auxiliary prime products and their dyadic lengths\\
\(H_i\), \(\mathscr A\), \(\mathcal E\)&Residue subgroups, a prime coset, and forbidden prime sets\\
\(B_0\), \(S\)&Finite exceptional-prime sets in the ordinary-indicator and primitive-form applications\\
\(\Delta_i=f_i^2D\)&Order discriminants sharing one fixed fundamental discriminant \(D\); here \(f_i\) denotes a conductor\\
\(P_Q,\Lambda_Q\), \(G\)&Primitive numerical values, their local shadow, and the applicable class group\\\hline
\end{tabular}
\end{center}
The modulus and the arithmetic data stay fixed. Definitions in each
application specify its groups and exceptional primes; \(h\) is
reserved for the additive shift throughout.
